\documentclass[10pt,letterpaper]{article}
\usepackage[T1]{fontenc}
\usepackage[utf8]{inputenc}
\usepackage{lmodern}
\usepackage[margin=0.76in]{geometry}
\usepackage[table]{xcolor}
\usepackage{array,tabularx,longtable,booktabs}
\usepackage{enumitem}
\usepackage{amssymb,amsmath}
\usepackage{hyperref}
\usepackage{fancyhdr,lastpage}
\usepackage{microtype}
\usepackage{parskip}
\definecolor{cdblue}{HTML}{203957}
\definecolor{cdgray}{HTML}{586677}
\hypersetup{colorlinks=true,linkcolor=cdblue,urlcolor=cdblue,pdftitle={Chapter 14: Advanced Version Control}}
\pagestyle{fancy}
\fancyhf{}
\lhead{\small\color{cdgray}Continuous Delivery Engineering}
\rhead{\small\color{cdgray}Study guide}
\cfoot{\thepage\ / \pageref{LastPage}}
\setlength{\headheight}{14pt}
\setlength{\emergencystretch}{2em}
\setlist[itemize]{leftmargin=*,itemsep=3pt,topsep=4pt}
\setlist[enumerate]{leftmargin=*,itemsep=3pt,topsep=4pt}
\begin{document}
\begin{center}
{\LARGE\bfseries\color{cdblue}Chapter 14: Advanced Version Control}\\[7pt]
{\large III: The Delivery Ecosystem}\\[6pt]
{\small Reference edition: 2026-10-30 $\bullet$ Original educational and operational material}
\end{center}
\noindent\color{cdblue}\rule{\linewidth}{1.2pt}\color{black}
\vspace{0.4em}
\section{Chapter scope and provenance}
Primary reading: Humble and Farley, Continuous Delivery (2010), Chapter 14, "Advanced Version Control"; Part III: The Delivery Ecosystem. This guide is an original synthesis and not a reproduction of the chapter. The book provides the foundational delivery principles; current AWS, Terraform and GitHub Actions references in the applied sections are supplementary.


\section{Learning outcomes}
\begin{enumerate}
\item Assess branch lifetime and integration delay.
\item Use release branches only for supported lifecycle needs.
\item Choose merge, revert, and feature toggle controls.
\end{enumerate}

\section{Core conceptual model}
Focus areas: branching and merging; mainline development; release branches.

\begin{itemize}
\item Input: a version-controlled change, explicit requirements, and a reproducible execution context.
\item Mechanism: controlled feedback, automation, evidence capture, and a stop condition where risk exceeds tolerance.
\item Output: a verifiable release candidate or a diagnosed failure with an accountable owner.
\end{itemize}

\section{Reading and analysis protocol}
\begin{enumerate}
\item Read the chapter and record five propositions, distinguishing an author claim from an example or assumption.
\item Identify two failure modes, one recovery technique, and one measurable acceptance condition.
\item Trace each technique to the stages and controls in the deployment-pipeline-architecture reference.
\item State how the proposed practice changes for a regulated system with segregated duties.
\end{enumerate}

\section{Detailed technical study notes}
\subsection{Mainline as integration point}
Treat trunk as the continuous source of integration truth and shorten the divergence of candidate changes.

\subsection{Branch lifetime}
Longer-lived branches accumulate collision and drift risk; measure open time and late merge failures.

\subsection{Release stabilization}
Where release branches are justified, define patch flow, test scope, merge-back policy and end-of-support.

\subsection{Feature toggles}
Use controlled toggles or abstraction boundaries to separate unfinished work from release readiness, with cleanup and ownership.

\subsection{Merge and revert}
A revert is a new change with potential effects on later work. Apply tests and document residual risk.

\subsection{Version audit}
Tie source ancestry, tags and actual deployed artifact digests together; a tag alone is not release evidence.

\section{Domain-specific challenge questions}
\begin{enumerate}
\item For mainline as integration point, identify the precondition, expected evidence, adversarial failure case and measurable corrective action.
\item For branch lifetime, identify the precondition, expected evidence, adversarial failure case and measurable corrective action.
\item For release stabilization, identify the precondition, expected evidence, adversarial failure case and measurable corrective action.
\item For feature toggles, identify the precondition, expected evidence, adversarial failure case and measurable corrective action.
\end{enumerate}
\section{Guided exercise}
Run a simulated branch-by-abstraction migration and compare merge conflict rates.

\subsection{Required evidence}
\begin{itemize}
\item A version-controlled design or executable artifact.
\item A failure demonstration, negative test or counterexample.
\item An improvement statement with baseline, target, owner and verification method.
\end{itemize}


\section{Engineering scenario}
A hotfix exists on the release branch but not on main; design the propagation control.

\begin{itemize}
\item Describe the failure propagation path and the earliest feasible detection point.
\item Recommend a minimal remediation and its rollback or containment plan.
\item Define a measure that would demonstrate risk reduction.
\end{itemize}

\section{Knowledge checks}
\begin{enumerate}
\item Which condition is the gate, and which action is the response when it fails?
\item What evidence distinguishes successful automation from an assumed successful release?
\item Where might environments, data, dependencies or people invalidate the technique?
\item What feedback could be moved earlier without losing necessary confidence?
\item Which assumptions are book principles, and which are modern supplementary implementation choices?
\end{enumerate}

\section{Completion rubric}
Score each dimension from 0 (missing) to 4 (demonstrated and independently verifiable): conceptual accuracy; reproducibility; negative testing; security and safety; documentation and traceability. Pass threshold: 16/20 with no zero on reproducibility or safety.


\section{Prerequisites and onward route}
Chapters 1--4 and prior dependency stages. Next: Chapters 15 onward.



\end{document}
